\documentclass[11pt,a4paper]{article}
\usepackage[margin=25mm]{geometry}
\usepackage{amsmath,amssymb,amsthm,hyperref}
\setlength{\emergencystretch}{2em}
\newtheorem{lemma}{Lemma}
\newtheorem{theorem}[lemma]{Theorem}
\newcommand{\gen}[1]{\langle#1\rangle}
\newcommand{\Sym}{\operatorname{Sym}}
\title{A trivial internal factor by moving the split}
\author{Human Compression project: a component of the L019 recovery}
\date{}
\begin{document}\maketitle
All actions are actual permutation actions. A perfect split at $i$
has two omission orbits $A,B$, with the marks before $i$ supported
and transitive on $A$, and those after $i$ supported and transitive
on $B$. Its crossing mark has the form
\[
 \rho_i=(a\ b)\alpha\beta,
 \qquad a\in A,\quad b\in B,
\]
where $\alpha,\beta$ are supported on $A\setminus\{a\}$ and
$B\setminus\{b\}$. Call an internal perfect split \emph{usable}
if at least one of these two internal factors is trivial. This term
asserts only the support property, not the separate leaf-kernel
conditions needed for an extension.

We acquire the complete actual-wing and faithful-augmentation proofs,
the boundary and internal perfect-split theorems, and the two
inheritance lemmas. The faithful-augmentation interface is the following:
if $2r>n+2$, both internal factors are nontrivial, and the left wing
is $\Sym(A)$, then either its rank $i$ satisfies $2i>|A|+2$, or
\[
 |B|=5,\quad \operatorname{rank}(G_{>i})=2,\quad
 |A|=m\text{ is even},\quad i=m/2+1.                 \tag{1}
\]
The right wing in (1) has order ten. This conclusion uses a faithful
restriction proved by the actual intersection property and an exact
filter of the acquired primitive classification; it does not rely
on quotient IP. The same statement holds in the reversed order.

\begin{lemma}\label{move}
Suppose a fractured parent string C-marking has an internal perfect split at $i$,
its left wing is the natural $S_m$ of rank $i=m/2+1$,
and $m$ is even. Suppose its old left internal factor is nontrivial. Then either the parent has a usable split, or
it has a perfect split at an index $k$ with $1\le k\le i-2$.
Here the parent diagram is connected and the old attachment point
is an endpoint anchor of the wing.
\end{lemma}
\begin{proof}
The nontrivial old internal factor gives $m\ge4$.
The wing inherits the intransitive-omission hypothesis.
Its noncommutation diagram is connected. For $m\ge5$, this follows
from the commuting-normal-factor obstruction in $S_m$.
For $m=4$, the rank is three; a disconnected string would have an
isolated mark central in $S_4$, impossible since its centre is trivial.
The boundary theorem, since $2i=m+2$, supplies a perfect wing split
at some $k$.

If $k=i-1$, reversed adjacent inheritance gives a parent perfect
split with a trivial internal factor, using $2i>m-1$.
If $1\le k\le i-2$, nonadjacent inheritance gives the second
conclusion, preserving either trivial factor if present.
It remains $k=0$.

The complement of that singleton has degree $m-1$, rank
$i-1=m/2$, and is full symmetric by the acquired nonsymmetric-wing
bound: every nonsymmetric actual wing has rank at most half its
degree. That bound's proof is independent of its final parent rank
inequality. This complementary wing has an anchor at its first end,
coming from the zero split. It also contains the original attachment
anchor at its last end: that point is moved by $\rho_{i-1}$,
whereas the removed singleton is moved only by $\rho_0$.
Thus the two-anchor theorem applies, with
$2(i-1)=m>(m-1)$. Its omission hypothesis is inherited once more.
Its diagram is connected because it is $S_{m-1}$; for $m=4$
this is the nonabelian rank-two group $S_3$.

Let $j\in\{1,\ldots,i-1\}$ be its perfect split.
If $j=1$, adjacent inheritance to the full left wing applies since
$2(i-1)=m>m-2$ and gives a trivial near factor. This then inherits
nonadjacently to the parent, since $i\ge3$ and $1\le i-2$.
If $j\ge2$, it inherits nonadjacently to the left wing.
For $j=i-1$ use reversed adjacent inheritance through the old
parent cut; otherwise it lies in $\{1,\ldots,i-2\}$ and
inherits nonadjacently to the parent. These are precisely the
claimed alternatives.
\end{proof}

\begin{lemma}\label{opposite}
Suppose a fractured string C-marking has perfect splits at $2$
and $i=r-3$, and the three regions between them are $C,D,B$,
where $|C|=|B|=5$ and $|D|=m-5$ for even $m$.
Suppose its actual wings on $C\cup D$ and $D\cup B$ are both
full symmetric of rank $m/2+1$, so $i=m/2+1$.
Then either it has a usable split or it has an internal perfect
split with both omission orbits larger than five.
\end{lemma}
\begin{proof}
The two wing groups are actual symmetric groups on their displayed
supports. The intersection property therefore identifies
\[
 T=\gen{\rho_3,\ldots,\rho_{i-1}}=\Sym(D).
\]
In particular $T$ is transitive on $D$, with rank
$t=i-3=m/2-2$ and degree $d=m-5$, so $2t=d+1$.
Since $D$ is nonempty, $m\ge6$. The case $m=6$ would give
a nonidentity single mark acting on a singleton, which is impossible.
Thus $m\ge8$, $d\ge3$, and $t\ge2$.
The group is connected for $d\ge5$ by the symmetric normal-factor
argument, and for $d=3$ because it is the nonabelian rank-two $S_3$.

The cut at $2$ remains perfect inside the symmetric left wing:
its left action is transitive on $C$, its middle action $T$ is
transitive on $D$, and $\rho_2$ fixes $B$. Similarly the cut
at $i$ remains perfect inside the symmetric right wing. This
justifies the two successive wing-omission transfers, so every
maximal omission of $T$ is intransitive. The inner endpoints of
the two crossing edges give its two endpoint anchors: all middle
marks except $\rho_3$, respectively $\rho_{i-1}$, fix them
by the unique-crossing commutations, and transitivity forces those
end marks to move them. The two-anchor theorem supplies a perfect
split $j$ of $T$.

If $j=3$, adjacent inheritance through the cut at $2$ inside the
left wing applies since $2t=d+1>d-1$. It gives a trivial near
factor. This split then inherits to the full parent through the
cut at $i$, nonadjacently because $t\ge2$ gives $3\le i-2$.
If $j=i-1$, use the reversed argument inside the right wing, then
inherit through the cut at $2$.

Otherwise $4\le j\le i-2$. Inherit the split nonadjacently
first into the left wing through its cut at $2$, and then to the
parent through the cut at $i$. Its two sides consist of $C$ plus
the nonempty left middle side, and $B$ plus the nonempty right
middle side. Both therefore have size greater than five.
If either internal factor is trivial, it is usable already;
otherwise the second conclusion is exactly what remains.
\end{proof}

\begin{theorem}\label{usable}
A fractured string C-marking of the natural $S_n$, with
$n\ge5$ and $2r>n+2$, has a usable split.
Consequently every natural symmetric C-marking with $n\ge7$
in this rank range has such a split.
\end{theorem}
\begin{proof}
We use strong induction on the degree, retaining the actual fracture
hypothesis. For degrees five and six the acquired complete extremal
classification gives the ordinary simplex, except for the outer
image in degree six. The latter has transitive omissions and is
excluded by the present hypothesis. The ordinary adjacent-transposition
marking has an internal pure split. Degrees below five cannot satisfy
the rank inequality, by the acquired bound $r\le n-1$.
For $n\ge7$, the internal-split theorem supplies an internal perfect
split $i$. If either internal factor is trivial, it is usable.
Assume both are nontrivial.

The actual-wing theorem makes at least one wing symmetric; reverse
the whole order to call it the left wing on $A$. If its rank
satisfies $2i>|A|+2$, apply the inductive hypothesis to that
smaller actual wing. It inherits fracture and IP, and its internal
usable split has index in $\{1,\ldots,i-2\}$.
Nonadjacent inheritance transfers it to the parent and preserves its
trivial factor. The only remaining case is (1).

Apply Lemma~\ref{move} to this critical even wing. We either have
a usable split, or obtain an internal perfect split $k$ with
$1\le k\le i-2$. Its left side is a nonempty set $C\subset A$,
and its right side is $D\cup B$, with $D=A\setminus C$ nonempty.
These are the actual three regions: the $k$-edge lies inside $A$,
and nonadjacent inheritance adds the whole old side $B$ to its right.
If an internal factor at $k$ is trivial, stop.

Otherwise choose a symmetric wing at this new parent split and
apply faithful augmentation again. If it has the strict smaller-degree
rank, induction and nonadjacent inheritance finish as before. If the
symmetric wing were on $C$, its opposite side would have size
$|D|+5>5$, so the exceptional alternative (1) is impossible.
The only new exceptional alternative consequently has its symmetric
wing on $D\cup B$, and its opposite side must be $C$ of size five
and rank two. Thus $k=2$; the new symmetric wing has degree
$|D|+5=m$, while the old one on $C\cup D$ had the same degree.
We have exactly the two opposite ends of Lemma~\ref{opposite}.

That lemma either gives a usable split or produces an internal perfect
split with both sides greater than five. If both its internal factors
are nontrivial, the actual-wing theorem and faithful augmentation now
force a strict smaller symmetric wing: neither opposite side has the
exceptional size five. Induction on that smaller degree and the same
nonadjacent inheritance give a usable split of the parent. All recursive
uses decrease the actual symmetric degree; moving a cut and treating the
two opposite ends are finite steps within the present induction case.

Finally, for $n\ge7$ the acquired initial-fracture theorem supplies
the fracture hypothesis directly from the natural symmetric marking
and its strict rank. This proves the last assertion.
\end{proof}

\paragraph{Attribution and outstanding requirements.}
The conclusion recovers the trivial-factor step in Section~8 of
Cameron--Fernandes--Leemans, \emph{The number of string C-groups
of high rank}, arXiv:2212.12723v2. Its smaller-wing strategy and
the historical rank classifications remain attributed. The new proof
uses the faithful augmentation already established in this recovery,
then moves a perfect split inside the critical even wing. The sole
failure of strict descent becomes two actual symmetric wings with a
symmetric middle intersection, to which the endpoint-anchor theorem
applies. No unsupported converse to a central twist, projected IP,
or suffix-centralizer assertion is used.

The complete endpoint-anchor manuscript and its distinct review are
acquired premises. A usable split as defined here is not
asserted to satisfy all leaf-kernel conditions. Those conditions, the
reversible counting correspondence and its exact base counts remain
separate obligations; the whole original is not yet credited.
\end{document}
